% Wu & Tu's abstract is a single paragraph of about 150 words, and it ends on the "we hope this
% informs transformers" note rather than on a number. Match both. Written last; the version
% below fixes the claim so the rest of the paper has a target.

The probabilistic transformer of \citet{wu2023probabilistic} models contextual word
representation as posterior inference in a conditional random field over dependency structure,
and its mean-field inference procedure turns out to resemble a transformer encoder. That model
is an encoder only. In this work we derive its autoregressive counterpart. We show that
causality, a single global energy function and contextuality cannot hold simultaneously, and
resolve the conflict with a directed chain of conditional random fields in which the prefix
enters each step as a frozen condition. We further obtain a language-model head without adding
any parameter outside the factor graph, by promoting the word variable from observed to latent;
tied input and output embeddings then follow as a consequence rather than a design choice.
Because the per-position subgraph is a tree, the vocabulary readout is available in closed form.
Experiments on \TODO{corpora} compare our model against a standard causal transformer and
against a looped transformer, the latter isolating the effect of syntactic structure from that
of weight sharing. \TODO{one sentence of result} We hope that this work helps extend the
probabilistic account of contextual representation to generation.
